\documentclass[11pt]{article}
\usepackage[margin=2.2cm]{geometry}
\usepackage{booktabs}
\usepackage{amsmath}
\usepackage{graphicx}
\usepackage{tabularx}
\usepackage{float}
\usepackage{enumitem}
\usepackage[hidelinks]{hyperref}
\setlist{itemsep=2pt}
\newcolumntype{Y}{>{\raggedright\arraybackslash}X}
\renewcommand{\arraystretch}{1.15}

\title{DMMST Version 2 --- explained from zero}
\author{Parsa}
\date{\today}

\begin{document}
\maketitle

\noindent This document explains, without assuming any background, what survival analysis is, what
our model does, what every loss and dataset is, what we ran in version 2, and what came out. Numbers
are from notebooks 01--06 (the language-model notebook 07 is not included).

\tableofcontents
\newpage

%==========================================================================
\section{The problem}

\subsection{What is survival analysis?}
We want to predict \textbf{when} something will happen to a patient --- death, relapse of a cancer,
recovery after a transplant. That ``something'' is called an \textbf{event}, and the time until it
happens is the \textbf{event time}.

It is not a normal prediction problem because of \textbf{censoring}: for many patients we never see
the event. A study runs for, say, 10 years; a patient who is still alive at the end, or who moved away
after 3 years, has no event time. We only know ``the event had not happened by time $T$''. Such a
patient is \textbf{censored} at time $T$. Throwing censored patients away would bias everything (they
are often the healthiest ones), so survival models are built to use the partial information they carry.

\subsection{What does a survival model output?}
Instead of one number, a survival model usually outputs a \textbf{survival curve} for each patient:
\[
  S(t) = \text{the probability that the event has \emph{not} happened by time } t.
\]
$S$ starts at 1 (at time 0 nothing has happened) and goes down. A high-risk patient's curve drops
fast, a low-risk patient's curve drops slowly. From the curve one can read a risk at any time, or an
expected event time.

\subsection{One event or several events?}
\begin{itemize}
  \item \textbf{Single event}: one kind of event, e.g.\ death.
  \item \textbf{Multiple events}: several kinds per patient, e.g.\ after a bone-marrow transplant:
        recovery, an adverse event, relapse, death. They can be
        \begin{itemize}
          \item \textbf{competing}: only the first one is seen (if you die of A you cannot later die of B);
          \item \textbf{semi-competing / non-exclusive}: several can happen to the same patient, in some
                order (recovery, then relapse, then death).
        \end{itemize}
\end{itemize}
Our paper (DMMST) is about a model that handles \textbf{multiple events}.

%==========================================================================
\section{Our model}

\subsection{The idea in one paragraph}
Each patient is a list of features (age, lab values, treatment, \dots). We turn every feature into a
small vector (a \textbf{token}), feed all tokens of a patient into a \textbf{transformer} (the same
kind of network used in language models), and get one vector that summarises the patient. On top of
that vector sit \textbf{heads} --- small networks that produce the outputs:
\begin{itemize}
  \item the \textbf{survival head} outputs one survival curve per event;
  \item the \textbf{regression head} (optional) outputs one predicted event time per event.
\end{itemize}

\subsection{The numeric embedding (paper \S2.2)}
How do you give a number like ``age = 63'' to a transformer? Two options:
\begin{itemize}
  \item \textbf{Discretise} it: cut age into bins (e.g.\ 60--65) and treat the bin as a word. Simple,
        but loses precision.
  \item \textbf{Numeric embedding} (what the paper proposes): every feature has its own learned vector,
        and that vector is \textbf{multiplied by the (standardised) value}. Age 63 and age 64 then give
        almost the same input, as they should.
\end{itemize}
We also tried \textbf{quantile} scaling (the value is replaced by its rank among all patients before
the multiplication).

\subsection{Representation: how to get one patient vector}
The transformer outputs a vector per token per layer. We tried different ways to combine them: over
\textbf{layers} (sum, mean, concatenate, or a special [CLS] token), over \textbf{tokens} (max or mean),
and using \textbf{all layers} or only the \textbf{last} one.

\subsection{How the survival head works ($L_{PCH}$)}
Time is cut into a few intervals (e.g.\ 0--14 days, 14--18 days, \dots). In each interval the head
predicts a \textbf{hazard}: how likely the event is in that interval if it has not happened yet.
Multiplying the ``no event'' chances interval by interval gives the survival curve. This is called a
\textbf{piecewise-constant hazard (PCH)} model.

The model is trained with $L_{PCH}$, the likelihood of what we actually observed:
\begin{itemize}
  \item for a patient with an event at time $T$: ``survive until $T$, then the event happens at $T$'';
  \item for a patient censored at time $T$: ``survive until $T$'' (nothing more is claimed).
\end{itemize}
This is how censored patients are used correctly.

%==========================================================================
\section{The losses}
A \textbf{loss} is the number the model tries to make small during training. $L_{PCH}$ alone already
gives a working model. The paper adds more losses; testing whether they help is a large part of
version 2.

\begin{table}[H]
\centering
\begin{tabularx}{\textwidth}{llY}
\toprule
Loss & Paper & What it asks of the model, in plain words \\
\midrule
$L_{PCH}$ & Eq.~3 & Make the observed event times (and censoring times) likely. The base loss. \\
$L_{rank}$ & Eq.~4 & \textbf{Between patients}: if patient A had the event before patient B, A should have
  been predicted as riskier than B (for the same event). \\
$L_{mul}$ & Eq.~5 & \textbf{Within one patient}: if a patient had event 1 before event 2, event 1 should
  have been predicted as riskier at that time. Needs $\geq 2$ events. \\
$L_{MAE}$ & Eq.~6 & For the regression head: the predicted time should be close to the real time. \\
$L_{MM}$ & Eq.~7 & For the regression head: the predicted times of the different events should agree
  with each other (``mismatch'' penalty). Needs $\geq 2$ events. \\
\bottomrule
\end{tabularx}
\end{table}

\paragraph{Best guess (Eq.~6).} For a censored patient we do not know the true time, only that it is
later than $T$. The ``best guess'' is an estimate of the real time made from all training patients
(Kaplan--Meier curve): ``patients who were still event-free at $T$ had their event on average at time
$X$''. $L_{MAE}$ can either ignore censored patients (R1) or use their best guess (R2).

\paragraph{Weight and $\sigma$.} $L_{rank}$ and $L_{mul}$ have a strength (weight) and a sharpness
($\sigma$). These were tuned. The paper writes Eqs.~4--5 with the wrong sign; we use the correct
direction (as in DeepHit) and a test checks it.

\subsection{The recipes we tested}
\begin{table}[H]
\centering
\begin{tabular}{ll@{\hspace{2cm}}ll}
\toprule
\multicolumn{2}{l}{Survival head} & \multicolumn{2}{l}{Regression head} \\
\midrule
S1 & $L_{PCH}$ only & R1 & $L_{MAE}$, observed events only \\
S2 & $L_{PCH} + L_{rank}$ & R2 & $L_{MAE}$ with the best guess \\
S3 & $L_{PCH} + L_{mul}$ (multi-event) & R3 & R2 + $L_{MM}$, observed time as reference \\
S4 & $L_{PCH} + L_{rank} + L_{mul}$ (multi-event) & R4 & R2 + $L_{MM}$, best guess as reference \\
\bottomrule
\end{tabular}
\end{table}
``---'' in a table means the recipe needs $\geq 2$ events and does not exist on that dataset.

We also removed one loss that was in earlier versions, \textbf{MMV} (from another paper, UniSurv): it
is not in our paper and does the same job as our regression head.

%==========================================================================
\section{The data}
\begin{table}[H]
\centering
\begin{tabularx}{\textwidth}{lllp{3.3cm}Y}
\toprule
Dataset & Real? & Patients & Events & What it is \\
\midrule
METABRIC & real & 1,904 & 1: death & breast-cancer patients; 9 features (genes + clinical) \\
SUPPORT & real & 8,873 & 1: death & seriously ill hospital patients; 14 features \\
EBMT & real & 2,279 & 4: recovery, adverse event, relapse, death & bone-marrow transplant patients;
  only 4 categorical features \\
hsa synthetic & simulated & 5,000 & 2, both can happen & from Tjandra et al.\ (2021); about 70\%
  censored at the end of the study \\
DeepHit synthetic & simulated & 30,000 & 2, competing & from the DeepHit paper (Lee et al.\ 2018) \\
\bottomrule
\end{tabularx}
\end{table}
METABRIC and SUPPORT are the standard datasets of SurvTRACE, our closest competitor, so we can compare
with its published numbers. \textbf{SEER} was planned as a large real dataset, but because of the
updated SEER+ data-access policies we could not get access to it.

%==========================================================================
\section{How we measured (the protocol)}

\subsection{Splits}
Each dataset is split at random into \textbf{60\% training} (the model learns from it), \textbf{10\%
validation} (used to choose settings and when to stop training) and \textbf{30\% test} (used only for
the final numbers). This is done \textbf{10 times} with different random splits, and every model gets
the same 10 splits. Results are written as \textbf{mean (standard deviation)} over the 10 splits ---
e.g.\ 0.667 (0.013). The standard deviation shows how much the result moves from split to split.

\subsection{Tuning (notebook 01)}
Every model has settings (learning rate, network size, \dots). We tried many combinations on the
\textbf{validation} split and kept the best --- for every model, not only ours, so the comparison is
fair. Then, with the chosen network, we tuned the loss weights ($L_{rank}$, $L_{mul}$, $L_{MM}$).

\subsection{Early stopping and ``patience''}
During training the model is checked on validation regularly. If it has not improved for a number of
checks (the \textbf{patience}), training stops. We use 30 for our loss and regression experiments
(Section~\ref{sec:fixes}).

\subsection{The scores}
\begin{table}[H]
\centering
\begin{tabularx}{\textwidth}{lYl}
\toprule
Score & Meaning & Good value \\
\midrule
$C_{td}$ & Of all pairs of patients, how often is the one who had the event first predicted as riskier?
  Measured at the 25\%, 50\% and 75\% points of the event times, then averaged. & higher; 0.5 = coin flip \\
Brier score & Squared error of the predicted probability ``event by time $t$'' versus what happened. & lower \\
IBS & Brier score averaged over the whole follow-up. & lower \\
IBLL & Like IBS, but with log-loss. & lower \\
AUC & Another standard way to measure ranking at a time point. & higher \\
Antolini $C$ & A $C_{td}$ version that uses the whole curve. & higher \\
D-calibration $p$ & Tests whether predicted probabilities match observed frequencies. Large test sets make it
  very strict. & higher ($>0.05$ = no problem found) \\
MAE & Average error of a predicted event time, in the dataset's time unit. Censored patients are handled
  with the best guess (``margin'') or ``pseudo-observations''. & lower \\
Harrell $C$ & Like $C_{td}$ but for a predicted \emph{time}: are earlier events predicted earlier? & higher \\
\bottomrule
\end{tabularx}
\end{table}
In short: \textbf{$C_{td}$ = does the model put patients in the right order; IBS = are its
probabilities right.}

%==========================================================================
\section{The models we compared against}
\begin{table}[H]
\centering
\begin{tabularx}{\textwidth}{lY}
\toprule
Model & What it is \\
\midrule
SurvTRACE & A transformer survival model (2022), run with the authors' own code. Our closest competitor. \\
DeepHit & Neural network that predicts the event-time distribution directly (2018). \\
DSM & Deep Survival Machines: a mixture of simple distributions (2021). \\
MENSA & A recent multi-event model (2024). \\
RSF & Random Survival Forest: many decision trees. A strong classic. \\
PC-Hazard & A simple neural network with the same piecewise hazard as our head. \\
DeepSurv & A neural version of the Cox model. \\
Cox & The classic statistical survival model (linear). \\
\bottomrule
\end{tabularx}
\end{table}
DeepHit, DSM and MENSA use our transformer with their own head; the others are separate models. On
multi-event data, single-event models are trained once per event.

%==========================================================================
\section{What each notebook did}
\begin{table}[H]
\centering
\begin{tabularx}{\textwidth}{lYr}
\toprule
Notebook & What it did & Runs \\
\midrule
01 Tuning & chose the settings of every model and the loss weights, on validation only & 576 \\
02 Tabular benchmark & all 9 models on METABRIC and SUPPORT (single event) & 180 \\
03 Multi-event benchmark & all 9 models on EBMT, hsa synthetic, DeepHit synthetic & 270 \\
04 Encoding / representation & numeric embedding vs quantile vs bins; 13 ways to pool the transformer & 320 \\
05 Survival losses & S1--S4 on all datasets & 160 \\
06 Regression head & R1--R4 alone, and every S $\times$ R together with the survival head (``joint'') & 720 \\
\bottomrule
\end{tabularx}
\end{table}

%==========================================================================
\section{The results, in plain words}

\subsection{Are we better than the other models? (02, 03)}
\begin{table}[H]
\centering
\resizebox{\textwidth}{!}{\begin{tabular}{lccccc}
\toprule
$C_{td}$ & METABRIC & SUPPORT & EBMT & hsa synth. & DeepHit synth. \\
\midrule
Best other model & SurvTRACE 0.678 & SurvTRACE 0.638 & RSF 0.549 & SurvTRACE 0.868 & DeepSurv 0.762 \\
Our best recipe & 0.672 & 0.620 & 0.537 & \textbf{0.868} & 0.744 \\
Our rank (of 9, S1) & 5 & 4 & 6 & 3 & 5 \\
\bottomrule
\end{tabular}}
\end{table}
\textbf{Plainly: our model is good but not the best.} It is always in the better half and ties
SurvTRACE on hsa synthetic, but on the other four datasets something else is slightly better (by
0.006--0.018). Our \textbf{calibration} (IBS) is good --- close to the best on four datasets, and
better than SurvTRACE on METABRIC. A sanity check also passed: our run of SurvTRACE gets almost its
published numbers (0.678 vs 0.691, 0.638 vs 0.640), so the setup is right.

\textbf{EBMT is hard for every model} (all between 0.50 and 0.55): its 4 features say little about the
outcome.

\subsection{Do the extra losses help? (05)}
\begin{table}[H]
\centering
\begin{tabular}{lccccc}
\toprule
$C_{td}$ & METABRIC & SUPPORT & EBMT & hsa synth. & DeepHit synth. \\
\midrule
S1: $L_{PCH}$ & 0.667 & 0.615 & 0.529 & 0.849 & 0.743 \\
S2: $+L_{rank}$ & 0.670 & 0.620 & 0.535 & 0.868 & 0.743 \\
S3: $+L_{mul}$ & --- & --- & 0.530 & 0.859 & 0.742 \\
S4: $+$both & --- & --- & 0.537 & 0.867 & 0.742 \\
\bottomrule
\end{tabular}
\end{table}
\begin{itemize}
  \item \textbf{$L_{rank}$ helps a little, on every dataset.} The gains are small; statistically clear
        only on hsa synthetic.
  \item \textbf{Their real benefit is reliability.} With $L_{PCH}$ alone, training sometimes goes wrong
        on one split: on hsa one split reached only 0.71 (the others about 0.86), and on EBMT several
        splits had very bad probabilities (Brier up to 0.6). With $L_{rank}+L_{mul}$ (S4) this did not
        happen: EBMT's IBS went from 0.338 to 0.233 and became much more stable.
  \item On DeepHit synthetic nothing changes --- $L_{PCH}$ alone is enough there.
\end{itemize}

\subsection{Does the numeric embedding help? (04)}
\begin{itemize}
  \item METABRIC: yes --- 0.666 versus 0.642 with bins.
  \item SUPPORT: no --- bins are slightly better (0.627 versus 0.617).
  \item How the transformer outputs are combined hardly matters (all within $\pm0.01$), except that
        the [CLS] token is clearly worst.
\end{itemize}

\subsection{Does the regression head work? (06)}
\begin{itemize}
  \item \textbf{R2 (best guess) is better than R1 (observed only)}: it orders the predicted times
        better on every dataset (Harrell $C$, e.g.\ 0.57 $\rightarrow$ 0.66 on hsa) and has a smaller
        error on 4 of 5 datasets.
  \item \textbf{$L_{MM}$ (R3/R4)} changes little.
  \item \textbf{Training both heads together} (``joint''):
        \begin{itemize}
          \item does not hurt the survival predictions, as long as $L_{rank}$ is used (on hsa:
                0.84--0.85 with $L_{rank}$, 0.72--0.82 without);
          \item on hsa it even makes the predicted times better (Harrell $C$ 0.69 on average, up to
                0.79, versus 0.66 alone);
          \item on DeepHit synthetic the predicted times get worse (0.55 versus 0.70), because joint
                models pick their best checkpoint by the survival score, not by the time prediction.
        \end{itemize}
\end{itemize}

\subsection{What this means for the paper}
\begin{table}[H]
\centering
\begin{tabularx}{\textwidth}{lY}
\toprule
Claim & Supported? \\
\midrule
``Better accuracy than other models'' & \textbf{No} --- competitive, upper half \\
``Well calibrated'' & Mostly (not on EBMT) \\
``$L_{rank}$ and $L_{mul}$ improve the model'' & Yes, but as \textbf{stability and calibration};
  accuracy gains are small \\
``The numeric embedding is better than bins'' & On METABRIC yes, on SUPPORT no \\
``The regression head works; the best-guess loss helps'' & Yes \\
``Training both heads together helps'' & On hsa yes; neutral elsewhere; worse time prediction on
  DeepHit synthetic \\
\bottomrule
\end{tabularx}
\end{table}

%==========================================================================
\section{Problems we found and fixed along the way}\label{sec:fixes}
Version 2 was run on Kaggle; several issues were found from the results and fixed before the final
numbers. They are listed so you know the final results are trustworthy.

\begin{table}[H]
\centering
\small
\begin{tabularx}{\textwidth}{lYY}
\toprule
Problem & What happened & Fix \\
\midrule
SurvTRACE crashed & its 2021 code used a NumPy name removed in NumPy 2 & restored the name \\
Ranking loss gave NaN & on whole-number times a formula ran outside its interval and overflowed & kept it
  inside the interval \\
Qwen (LLM) crashed & a library version clash on Kaggle & bypassed the check \\
IBS/IBLL/AUC missing & some test patients were followed longer than any training patient & handle them
  correctly (cap the time) \\
Brier wrong on hsa & 70\% of hsa patients (censored at the end) were silently dropped from the Brier
  score of the neural models & keep them \\
Joint models collapsed & the regression loss measured in days (hundreds) drowned the survival loss &
  measure time as a fraction of follow-up \\
Some joint runs never learned & they start slowly and were stopped too early (patience 10) & patience 30
  (also for notebook 05, to be fair) \\
``Last layer'' runs crashed & wrong layer index (notebook 04) & use the real last-layer number \\
Kaggle killed long sessions & runs went past the 12-hour limit and nothing was saved & notebooks stop at
  11\,h\,20 and save \\
\bottomrule
\end{tabularx}
\end{table}

%==========================================================================
\section{The code, briefly}
The code lives in the \texttt{dmmst} repository (a fork of the open-source \textbf{SAT} library,
extended by us). You never need to call most of it directly: the notebooks do. Here is what each part
is for.

\subsection{How one run works}
Every experiment is a short pipeline of Python programs, each started with a \textbf{configuration}:
\begin{center}
\texttt{prepare\_data} $\rightarrow$ \texttt{train\_tokenizer} $\rightarrow$
\texttt{train\_labeltransform} $\rightarrow$ \texttt{finetune}
\end{center}
\begin{itemize}
  \item \texttt{prepare\_data} reads the data; \texttt{train\_tokenizer} turns features into tokens;
  \item \texttt{train\_labeltransform} makes the time intervals and, per split, the best guesses;
  \item \texttt{finetune} trains the model, scores it and saves \texttt{metrics.json}.
\end{itemize}
The first three run once per dataset (the label step once per split); \texttt{finetune} runs once per
model, recipe and split. Baselines use \texttt{baselines.py} or \texttt{survtrace\_official.py}
instead of \texttt{finetune}.

\subsection{Configuration (\texttt{conf/}, Hydra)}
Settings are not written in the code but in YAML files, combined by the \textbf{Hydra} library. A run
is described by a few overrides on the command line, e.g.
\begin{center}
\small\texttt{python -m sat.finetune experiments=multievent/survival dataset=hsa\_synthetic\_me}\\
\small\texttt{tasks=v2\_survival v2\_rank\_coeff=1.0 seed=3 split\_seed=3}
\end{center}
\begin{table}[H]
\centering
\small
\begin{tabularx}{\textwidth}{lY}
\toprule
Folder / file & What it holds \\
\midrule
\texttt{conf/experiments/} & one group per dataset family; \texttt{defaults.yaml} holds all shared
  settings (\texttt{v2\_rank\_coeff}, \texttt{v2\_mul\_coeff}, \texttt{v2\_l1\_type}, \dots) \\
\texttt{conf/tasks/v2\_*.yaml} & which heads the model has (survival, regression, or both) \\
\texttt{conf/tasks/losses/v2.yaml} & the loss recipe: which losses are on and their weights \\
\texttt{conf/tasks/metrics/v2.yaml} & the metrics computed at evaluation \\
\bottomrule
\end{tabularx}
\end{table}

\subsection{The library (\texttt{sat/})}
\begin{table}[H]
\centering
\small
\begin{tabularx}{\textwidth}{lY}
\toprule
Path & What it does \\
\midrule
\texttt{prepare\_data.py}, \texttt{data/} & read each dataset into one common format;
  \texttt{data/splitter.py} makes the exact 60/10/30 splits from a seed \\
\texttt{train\_labeltransform.py} & time intervals and, \textbf{from the training split only}, the
  Kaplan--Meier best guesses and censoring weights \\
\texttt{models/heads/embeddings.py} & the numeric embedding (feature vector $\times$ value) \\
\texttt{models/heads/mtl.py} & the whole model: transformer + the configured heads \\
\texttt{models/heads/survival.py}, \texttt{regression.py} & survival head (hazards $\rightarrow$
  curve) and regression head (predicted time) \\
\texttt{loss/survival/nllpchazard.py} & $L_{PCH}$ \\
\texttt{loss/ranking/sample.py}, \texttt{multievent.py} & $L_{rank}$ (between patients) and $L_{mul}$
  (between events of one patient) \\
\texttt{loss/regression/l1.py}, \texttt{mismatch.py} & $L_{MAE}$ (R1/R2) and $L_{MM}$ (R3/R4) \\
\texttt{loss/factory.py}, \texttt{meta.py} & build the recipe: only losses with weight $>0$, added up \\
\texttt{finetune.py} & trains (Hugging Face Trainer), early-stops on validation $C_{td}$, scores
  validation and test, writes \texttt{metrics.json} \\
\texttt{evaluate/survtrace\_metrics.py} & $C_{td}$ and Brier at the 25/50/75\% time points \\
\texttt{evaluate/v2\_metrics.py} & IBS, IBLL, AUC, Antolini, calibration, MAE, Harrell $C$ \\
\texttt{evaluate/shared.py} & the same scoring for the baselines \\
\texttt{baselines.py} & Cox, RSF, DeepSurv, PC-Hazard (scikit-survival / pycox) \\
\texttt{survtrace\_official.py} & the authors' SurvTRACE code on our splits \\
\texttt{llm.py} & the language-model experiments (notebook 07) \\
\texttt{utils/km.py} & Kaplan--Meier ``best guess'' \\
\bottomrule
\end{tabularx}
\end{table}

\subsection{Running many experiments (\texttt{scripts/})}
\begin{table}[H]
\centering
\small
\begin{tabularx}{\textwidth}{lY}
\toprule
File & What it does \\
\midrule
\texttt{runner.py} & the \textbf{Runner}: launches runs in parallel, skips finished ones, retries
  failures once, stops safely before Kaggle's time limit, collects all results into one table,
  paired significance tests \\
\texttt{v2.py} & the \textbf{version-2 plan}: datasets, how each model is launched, search grids, loss
  recipes, reading \texttt{tuned.json}, tables \\
\texttt{verify\_protocol.py} & checks the protocol (split sizes, no test data in training files) \\
\texttt{prepare\_ebmt.py}, \texttt{prepare\_public\_synthetic.py}, \texttt{fetch\_survtrace.py} &
  download or build EBMT, the synthetic datasets and the SurvTRACE code \\
\bottomrule
\end{tabularx}
\end{table}

\subsection{The notebooks (\texttt{notebooks/version\_2/})}
The notebooks are \textbf{generated} by \texttt{build\_notebooks.py} from shared cells, so they all set
up the same way. On Kaggle each notebook (1) clones the code from GitHub, (2) installs packages,
(3) copies back finished runs from an attached previous output, (4) reads \texttt{tuned.json},
(5) builds its list of runs and gives it to the Runner, (6) writes tables and a zip. Each run leaves
two files: \texttt{<run>.json} (all metrics for validation and test) and \texttt{<run>.meta.json}
(exactly how it was run, including the code version).

\subsection{Tests (\texttt{tests/})}
159 automated tests check, among other things, that our metrics equal the reference libraries, that
the ranking losses push in the right direction, and the edge cases fixed during version 2. They run
with \texttt{pytest}.

%==========================================================================
\section{Where everything is}
\begin{itemize}
  \item Results of each notebook: \texttt{results/version\_2/<notebook>/}
  \item Figures and reports: \texttt{results/version\_2/analysis/}
        \begin{itemize}
          \item \texttt{REPORT.md} --- short comparison of our losses versus other models;
          \item \texttt{report\_latex/report.pdf} --- 3-page report;
          \item \texttt{report\_latex/full\_report.pdf} --- 16-page report with every table;
          \item \texttt{report\_latex/explained.pdf} and \texttt{EXPLAINED.md} --- this document.
        \end{itemize}
  \item The plan of version 2: \texttt{notebooks/version\_2/VERSION.md}
  \item Not yet included: notebook 07 (the language-model part of the paper, \S3) and sequential data
        (\S4.2, planned for version 3).
\end{itemize}

\end{document}
